% Generado por scripts/probe_computed_params.py -- no editar a mano.
\begin{table}[ht]
\centering
\caption{What the parameter tuple contributes, and how it is used. All rows share one BM25 index over the catalog text plus the parameter tokens ($k_1{=}0.60$, $b{=}0.35$) and differ only in how the query is formed, except the last, which uses the same extracted tuple as an exact filter instead of as query terms. Canonical test set, $n=16{,}590$.}
\label{tab:param_tokens}
\begin{tabular}{lcccc}
\toprule
\textbf{Query representation} & \textbf{Acc@1} & \textbf{Recall@5} & \textbf{MRR@10} & \textbf{Parent Acc@1} \\
\midrule
Query text alone & 0.871 & 0.934 & 0.898 & 0.973 \\
Parameter tuple alone & 0.219 & 0.760 & 0.431 & 0.219 \\
Query text + tuple read from the catalog record & 0.974 & 0.985 & 0.979 & 0.985 \\
Query text + tuple extracted from the query & 0.974 & 0.984 & 0.979 & 0.986 \\
\midrule
Same extracted tuple, as an exact filter & 0.903 & 0.967 & 0.917 & 0.985 \\
\bottomrule
\end{tabular}
\end{table}
